\paragraph{\trEnergy。}\label{sub:energy}
从 $L_0$、$L_1$ 与一个\trFloerDatum{}$(H,J)$ 出发。由于\trExact{}性，存在函数
$h_{L_k}\in C^\infty(L_k,\RR)$ 使 $dh_{L_k}=\theta|L_k$。在\trPathSpace{}%
$\{y\in C^\infty([0;1],M):y(0)\in L_0,\ y(1)\in L_1\}$ 上的
$H$-\trPerturbation\trActionFunctional{}为
\begin{equation}\label{eq:perturbed-action}
A_H(y)=\left(\int -y^*\theta+H(t,y(t))\,dt\right)+h_{L_1}(y(1))-h_{L_0}(y(0)).
\end{equation}
$A_H$ 的\trCriticalPoint{}集恰是 $\scr{C}(L_0,L_1)$，而\eqref{eq:floer-equation} 正
是某个 $L^2$ \trMetric{}下的负\trGradient{}流方程。因此，对任何
$u\in\scr{M}_Z(y_0,y_1)$ 都有\trApriori{}界
$E(u)=\int_Z|\partial u/\partial s|^2=A_H(y_0)-A_H(y_1)$。要把这一点推广到
\eqref{eq:inhomogeneous-equation}，人们对每个\trConnectedComponent{}$C\subset
\partial S$ 选取一个具有与前面同样性质的函数 $h_{L_C}\in C^\infty(L_C,\RR)$。于是对
$u\in\scr{M}_S(\{y_\zeta\})$，
\begin{equation}\label{eq:energy-definition}
E(u)\stackrel{\mathrm{def}}{\equiv}\int_S\tfrac12|Du-Y|^2
=\sum_{\zeta\in\Sigma^\pm}\mp A_{H_\zeta}(y_\zeta)-\int_S R(z,u(z)).
\end{equation}
每个 $h_{L_C}$ 的选取都带有歧义（相差一个常数；在不\trConnected{}的情形下则相差一个%
\trLocallyConstant{}函数），但对 $\zeta$ 求和时这一歧义会抵消，因此第一项是有意义
的。至于第二项，\trCurvature{}被积项 $R\in\Omega^2(S,\scr{H})$ 在 $S$ 的局部坐标
$(s,t)$ 下由
\begin{equation*}
R=\bigl(\partial_sK(\partial_t)-\partial_tK(\partial_s)
+\{K(\partial_s),K(\partial_t)\}\bigr)\,ds\wedge dt.
\end{equation*}
给出。这里 $\{\cdot,\cdot\}$ 是\trPoisson{}括号。由于 $K$ 在 $S$ 的诸端上为零，这一
积分\trBounded{}且界独立于映射 $u$，从而对 $u\in\scr{M}_S(\{y_\zeta\})$，$E(u)$ 亦然。
